% ====================================================================
% DATEI: content/11_operators_full_de.tex
% Vollständige Operatorfamilie F, G, H, Q, R, S, U, Psi, E_int
% Core 1.3.3 — kanonische Version
% ====================================================================

\section{Evolution und Interaktionsoperatoren}
\label{sec:operators-full}

Dieses Kapitel erweitert den in
Abschnitt~the formal model chapter eingeführten abstrakten Evolutionsoperator
zu der vollständigen Operatorfamilie, die in der Ontology of
Continua (OC) verwendet wird. Es fasst die Definitionen aus archival source corpus
zusammen und stellt sie in einer kompakten, vertikal konsistenten Form dar,
die für Core~1.3.3 geeignet ist.

Während dieses Abschnitts wird ein Kontinuum als
\[
  K(t) =
  \big(
    \Omega(t), A(t), P(t), J(t),
    \Theta(t), \partial\Omega(t), C(t), k(t)
  \big),
\]
mit den in Abschnitt~the formal model chapter definierten Komponenten geschrieben.
Der Einbettungsraum sei \(M\), mit Achsen \(A(M)\) und Beschränkungen,
die die zulässigen Zustände von \(K\) einschränken.

\subsection{Zerlegung des Evolutionsoperators}

Die strukturelle Evolution eines Kontinuums über einen infinitesimalen
Zeitschritt ist durch einen Operator beschrieben
\[
  E : K(t) \longrightarrow K(t+dt).
\]
Statt einen einzelnen monolithischen Operator anzugeben, zerlegt OC \(E\)
in eine Familie von Operatoren, die auf unterschiedliche Komponenten des
Kontinuums wirken:
\[
  E = (F, G, H, Q, R, S, U),
\]
wobei
\begin{itemize}
  \item \(F\) die Dynamik der Flüsse \(J\) steuert;
  \item \(G\) die Dynamik der Potentiale \(P\) steuert;
  \item \(H\) die Dynamik der Schwellwerte \(\Theta\) steuert;
  \item \(Q\) die Dynamik der Zyklen \(C\) steuert;
  \item \(R\) die Dynamik des Randes \(\partial\Omega\) steuert;
  \item \(S\) die strukturelle Rekonfiguration der Achsen \(A\) und
        der internen Konnektivität steuert;
  \item \(U\) die Evolution der Continuumness \(k(t)\) steuert.
\end{itemize}

Formal kann dies geschrieben werden als
\begin{align*}
  J(t+dt)              &= F\big(J(t), P(t), A(t), \Theta(t), M\big), \\
  P(t+dt)              &= G\big(P(t), A(t), J(t), \Theta(t), M\big), \\
  \Theta(t+dt)         &= H\big(\Theta(t), P(t), J(t), M\big), \\
  C(t+dt)              &= Q\big(C(t), P(t), J(t), \Theta(t)\big), \\
  \partial\Omega(t+dt) &= R\big(\partial\Omega(t), P(t), J(t), \Theta(t), M\big), \\
  A(t+dt)              &= S\big(A(t), P(t), J(t), \Theta(t), M\big), \\
  k(t+dt)              &= U\big(k(t), \Omega(t), A(t), P(t), J(t),
                               \Theta(t), C(t), \partial\Omega(t), M\big).
\end{align*}

Die Operatoren sind durch die Axiome von OC eingeschränkt, insbesondere die
Monotonie der Dimension, die Unmöglichkeit einer Evolution außerhalb des
Einbettungsraums sowie die Definition von Geburt und Tod.

\subsection{Operator \texorpdfstring{\texorpdfstring{$F$}{F}}{F}: Flussdynamik}

Der Flussoperator \(F\) steuert die zeitliche Änderung der Flüsse:
\[
  J(t+dt) = F\big(J(t), P(t), A(t), \Theta(t), M\big).
\]

Konzeptionell kodiert \(F\):
\begin{itemize}
  \item Erhaltungs- oder Bilanzgesetze (z.\,B. Kontinuitätsgleichungen,
        Kirchhoff-Summen, Erhaltung der Wahrscheinlichkeit);
  \item konstitutive Beziehungen, die Flüsse mit Potentialgradienten
        verknüpfen (z.\,B. Fick-, Ohm-, Fourier-Gesetze, Reaktionskinetik);
  \item die strukturelle Einteilung von Flüssen in stützende, kritische und
        destruierende Komponenten.
\end{itemize}

Eine generische Zerlegung lautet
\[
  J(t) = J_{\mathrm{support}}(t)
       + J_{\mathrm{critical}}(t)
       + J_{\mathrm{kill}}(t),
\]
wobei \(F\) diese Struktur erhält:
\begin{align*}
  J_{\mathrm{support}}(t+dt)
    &= F_{\mathrm{support}}\big(J, P, A, \Theta, M\big), \\
  J_{\mathrm{critical}}(t+dt)
    &= F_{\mathrm{critical}}\big(J, P, A, \Theta, M\big), \\
  J_{\mathrm{kill}}(t+dt)
    &= F_{\mathrm{kill}}\big(J, P, A, \Theta, M\big).
\end{align*}

Auf verschiedenen Ebenen \(K_x\) haben diese Klassen unterschiedliche
Ausprägungen:
\begin{itemize}
  \item in \(K_2\) (physikalische Kontinua) sind Flüsse diffusive, konvektive
        und feldvermittelte Ströme;
  \item in \(K_3\) (chemische Kontinua) entsprechen Flüsse Reaktions- und
        Transportraten;
  \item in \(K_4\)–\(K_5\) (protokelluläre und bioelektrische Kontinua) sind
        Flüsse Ionenströme, osmotische Flüsse und Stoffwechsel-Flüsse;
  \item in \(K_6\)–\(K_8\) (kognitive und soziale Kontinua) werden Flüsse zu
        Flüssen von Information, Aufmerksamkeit, Ressourcen und Verpflichtungen.
\end{itemize}

\subsection{Operator \texorpdfstring{\texorpdfstring{$G$}{G}}{G}: Potentialdynamik}

Der Potentialoperator \(G\) steuert die Änderung der Potentiale:
\[
  P(t+dt) = G\big(P(t), A(t), J(t), \Theta(t), M\big).
\]

In vielen konkreten Modellen wird \(G\) als Differentialrelation geschrieben
\[
  \frac{dP}{dt} = J,
\]
gegebenenfalls mit zusätzlichen nichtlinearen Termen für Dissipation,
Antrieb und Kopplung.
Auf struktureller Ebene verlangt OC nur:
\begin{itemize}
  \item Potentiale reagieren auf Flüsse und können diese wiederum einschränken;
  \item Änderungen der Potentiale respektieren die Vorzeichen- und
        Ungleichungsstruktur, die durch \(\Theta\) vorgegeben ist;
  \item bei Abwesenheit von Flüssen \(J = 0\) streben Potentiale gegen
        Schwellflächen oder relaxieren in intern bevorzugte Konfigurationen.
\end{itemize}

Unterschiedliche Klassen von Potentialen (energetisch, chemisch, elektrochemisch,
informational, normativ) sind Beispiele derselben strukturellen Rolle:
Einschränkungen und antreibende Kräfte auf \(\Omega(K)\) zu kodieren.

\subsection{Operator \texorpdfstring{\texorpdfstring{$H$}{H}}{H}: Schwellenwertdynamik}

Schwellwerte sind nicht statisch; sie können sich unter dem Einfluss von
Potentialen, Flüssen und Einbettungsrestriktionen anpassen, driften oder
rekonfigurieren. Der Schwellwertoperator \(H\) ist definiert durch
\[
  \Theta(t+dt) = H\big(\Theta(t), P(t), J(t), M\big).
\]

Beispiele sind:
\begin{itemize}
  \item Anpassung von Aktivierungskurven ionischer Kanäle im Nervengewebe;
  \item Plastizität der Existenzbereiche von Protokellen in wechselnden
        Umgebungen;
  \item Verschiebung von Normen und institutionellen Schwellwerten in
        sozialen Systemen;
  \item Änderungen von Kohärenz- oder Konsistenzschwellwerten in
        theoretischen Systemen.
\end{itemize}

Formal muss \(H\) erfüllen:
\begin{itemize}
  \item Kompatibilität mit den Einbettungsrestriktionen: Schwellwertänderungen
        dürfen keine Konfigurationen verlangen, die durch \(M\) verboten sind;
  \item Erhaltung der Taxonomie der Schwellwerte
        (\(\Theta_{\mathrm{exist}}, \Theta_{\mathrm{stab}},
        \Theta_{\mathrm{crit}}, \Theta_{\mathrm{dim}},
        \Theta_{\mathrm{death}}, \Theta_{\mathrm{expr}}, \Theta_{\mathrm{embed}}\));
  \item lokale Stetigkeit: kleine Änderungen in Potentialen oder Flüssen dürfen
        zu kleinen Schwellwertänderungen führen, außer an kritischen Punkten,
        an denen Bifurkationen erlaubt sind.
\end{itemize}

\subsection{Operator \texorpdfstring{\texorpdfstring{$Q$}{Q}}{Q}: Zyklische Dynamik}

Zyklen \(C(t)\) repräsentieren geschlossene Trajektorien, die die
Organisation des Kontinuums aufrechterhalten. Ihre Evolution wird gesteuert durch
\[
  C(t+dt) = Q\big(C(t), P(t), J(t), \Theta(t)\big).
\]

Strukturell behandelt \(Q\):
\begin{itemize}
  \item \emph{Geburt} neuer Zyklen, wenn sich Flüsse im Zustandsraum fern von
        \(\partial\Omega\) schließen;
  \item \emph{Stabilisierung} oder \emph{Verstärkung} vorhandener Zyklen,
        wenn stützende Flüsse dominieren;
  \item \emph{Abschwächung} und \emph{Zerstörung} von Zyklen bei destruierenden
        Flüssen oder Schwellwertverletzungen.
\end{itemize}

Zykluskomplexe sind besonders wichtig. Sei \(C_{\max}(t)\) die maximale Menge
paarweise kompatibler Zyklen, die das Kontinuum stützen. Dann besagt Satz~9 in
Abschnitt~the historical result-boundary chapter, dass der Tod mit \(C_{\max}(t) = \emptyset\)
zusammenfällt. Der Operator \(Q\) beteiligt sich daher direkt an
Leben--Tod-Übergängen.

\subsection{Operator \texorpdfstring{\texorpdfstring{$R$}{R}}{R}: Randdynamik}

Der Randoperator \(R\) steuert die Geometrie und Position des Randes
\(\partial\Omega\):
\[
  \partial\Omega(t+dt) =
    R\big(\partial\Omega(t), P(t), J(t), \Theta(t), M\big).
\]

Auf struktureller Ebene muss \(R\):
\begin{itemize}
  \item mit den Schwellwertdefinitionen konsistent sein:
        \(\partial\Omega = \{ s \mid \exists k: f_k(s) = 0 \}\);
  \item Expansion, Kontraktion und Bifurkation von \(\Omega\) zulassen;
  \item die Effekte von Flüssen und Schwellwertverschiebungen auf
        zugängliche Regionen kodieren.
\end{itemize}

In konkreten Implementierungen kann \(R\) viele Formen annehmen:
\begin{itemize}
  \item in \(K_2\) kann es der Bewegung von Phasengrenzen,
        Perkolationsschwellwerten oder Kohärenzoberflächen entsprechen;
  \item in \(K_3\)–\(K_4\) umfasst es die Dynamik von Membranen,
        Schnittstellen und Kompartiment-Grenzen;
  \item in \(K_5\) beschreibt es Schwellwerte der Erregbarkeit und
        refraktären Regionen im Leitfähigkeitsraum;
  \item in \(K_7\)–\(K_8\) kann es institutionelle und rechtliche
        Grenzen repräsentieren.
\end{itemize}

Abschnitt~the boundary geometry chapter enthält eine erweiterte Darstellung der
Randgeometrie und Patch-Modelle.

\subsection{Operator \texorpdfstring{\texorpdfstring{$S$}{S}}{S}: Strukturelle
Rekonfiguration}

Der Strukturoperator \(S\) steuert Änderungen an Achsen und interner
Konfiguration:
\[
  A(t+dt) = S\big(A(t), P(t), J(t), \Theta(t), M\big).
\]

Typische Rollen von \(S\) sind:
\begin{itemize}
  \item Aktivierung oder Deaktivierung von Achsen (z.\,B. Freischalten latenter
        Freiheitsgrade);
  \item Neuverkabelung von Konnektivitätsmustern innerhalb des Kontinuums;
  \item Grobkörnige Abbildung oder Verfeinerung effektiver Achsen unter
        Renormierung oder Abstraktion.
\end{itemize}

Wesentlich ist, dass \(S\) der Monotonie der Dimension unterliegt:
\[
  \dim A(t+dt) \ge \dim A(t).
\]
Wird eine neue Achse \(A_{\mathrm{new}}\) hinzugefügt, so dass
\(A_{\mathrm{new}} \notin \mathrm{span}_{\mathrm{str}}(A(t))\), tritt ein
Dimensionswechsel ein und das Kontinuum wird zu \(K_{x+1}\), anstatt auf demselben
Niveau zu bleiben. Dies wird im Folgenden durch die Geburtsoperatoren in
Abschnitt~the birth operators formalisiert discussion.

\subsection{Operator \texorpdfstring{\texorpdfstring{$U$}{U}}{U}: Continuumness-Dynamik}

Der Operator \(U\) steuert die Evolution der Continuumness \(k(t)\):
\[
  k(t+dt) =
  U\big(k(t), \Omega(t), A(t), P(t), J(t),
        \Theta(t), C(t), \partial\Omega(t), M\big).
\]

Mit der vereinheitlichten Definition von \(k(t)\) aus Abschnitt~the formal model chapter
bewertet \(U\), wie Änderungen im Zustandsraum, in Achsen, Flüssen,
Schwellwerten, Zyklen und Rand die Viabilität des Kontinuums beeinflussen.
Qualitativ gilt:
\begin{itemize}
  \item stützende Flüsse und stabile Zyklen erhöhen oder erhalten \(k(t)\);
  \item destruierende Flüsse und Zyklusverlust senken \(k(t)\);
  \item Annäherung an \(\partial\Omega\) reduziert \(k(t)\);
  \item Expansion von \(\Omega\) und Anreicherung von Zykluskomplexen erhöhen
        \(k(t)\).
\end{itemize}

Der Tod entspricht \(k(t) \to 0\) zusammen mit \(\Omega = \emptyset\).
Geburt entspricht dem Auftreten eines neuen Kontinuums mit \(k(t) > 0\) in einer
zuvor leeren Region des Möglichkeitsraums.

\subsection{Geburtsoperatoren \texorpdfstring{$\Psi_{x\to x+1}$}{Psi\_x→x+1}}
\label{sec:birth-operators}

Dimensionswechsel zwischen Ebenen werden durch Geburtsoperatoren vermittelt
\[
  \Psi_{x\to x+1} : (K_x, M_x) \longrightarrow (K_{x+1}, M_{x+1}).
\]

Strukturell ist \(\Psi_{x\to x+1}\) definiert, sobald die folgenden Bedingungen
erfüllt sind:
\begin{enumerate}
  \item \textbf{Neue Unterschiede.}\
        Es gibt eine Klasse von Unterschieden außerhalb des aktuellen
        strukturellen Spanns:
        \[
          D_{\mathrm{new}}\not\subseteq
          \mathrm{span}_{\mathrm{str}}(A(K_x)).
        \]
  \item \textbf{Spannung oberhalb der Dimensionsschwelle.}\
        Die strukturelle Spannung überschreitet die Dimensionsschwelle:
        \[
          T(K_x,t) > \Theta_{\mathrm{dim}}(K_x).
        \]
  \item \textbf{Verfügbare Achse im Einbettungsraum.}\
        Der Einbettungsraum \(M_x\) enthält mindestens eine Achse, die zur
        Aufnahme der neuen Unterschiede geeignet ist:
        \[
          A_{\mathrm{new}} \in A(M_x) \setminus A(K_x).
        \]
  \item \textbf{Nicht-leere zulässige Region.}\
        Es gibt eine nichtleere Menge von Zuständen \(\Omega(K_{x+1})\), die mit
        der neuen Achse und den Schwellwerten kompatibel ist.
\end{enumerate}

Die Wirkung von \(\Psi_{x\to x+1}\) lässt sich zusammenfassen als
\begin{align*}
  A(K_{x+1})      &= A(K_x) \cup \{A_{\mathrm{new}}\}, \\
  \Omega(K_{x+1}) &\subseteq
     \Omega(M_{x+1}) \cap
     \big\{ s \mid \Theta(K_{x+1})(s) \le 0 \big\}, \\
  M_{x+1}         &\supset M_x,
\end{align*}
mit allen anderen Komponenten von \(K_{x+1}\) (Flüsse, Potentiale, Zyklen,
Continuumness) durch das allgemeine strukturelle Gerüst bestimmt.

Geburtsoperatoren sind \emph{minimal} und \emph{irreversibel}:
jede nichttriviale Nutzung von \(A_{\mathrm{new}}\) konstituiert das neue
Kontinuum \(K_{x+1}\); eine Rückkehr zur vorherigen Dimensionalität würde die
Zerstörung von \(\Omega(K_{x+1})\) erfordern, was als Tod und nicht als
reversible Vereinfachung interpretiert wird.

\subsection{Interaktionsoperator \texorpdfstring{$E_{\mathrm{int}}$}{E\_int}}

Wenn zwei Kontinua \(K_a\) und \(K_b\) interagieren, wird deren gemeinsame
Dynamik durch einen Interaktionsoperator beschrieben
\[
  E_{\mathrm{int}} : (K_a, K_b, M)
    \longrightarrow (K_a', K_b', M'),
\]
wobei \(M'\) den aktualisierten Einbettungsraum für das kombinierte System
bezeichnet.

Auf struktureller Ebene modifiziert \(E_{\mathrm{int}}\) Potentiale, Flüsse,
Schwellwerte, Zyklen und gegebenenfalls Achsen jedes Kontinuums. Es werden
mehrere Regime unterschieden:

\begin{itemize}
  \item \textbf{Konkurrenz.}\
        Flüsse eines Kontinuums behindern die Aufrechterhaltung von Zyklen im
        anderen. Typischerweise können \(k_a\) und \(k_b\) nicht beide unbegrenzt
        wachsen.
  \item \textbf{Parasitismus.}\
        Ein Kontinuum entzieht einem anderen stützende Flüsse und erhöht dadurch
        die eigene Continuumness auf Kosten seines Wirts.
  \item \textbf{Symbiose.}\
        Gekoppelte Flüsse erhöhen die Continuumness beider Kontinua; gemeinsame
        Zyklen können entstehen, die das Paar stabilisieren.
  \item \textbf{Fusion.}\
        Achsen und Potentiale vereinen sich zu einem neuen Kontinuum
        \(K_{\mathrm{fusion}}\) mit zusammengeführtenden Zustandsraum und
        Schwellwerten. Formal entspricht die Fusion der Geburt eines
        höherdimensionalen Kontinuums über einen zusammengesetzten Geburtsoperator.
\end{itemize}

Interaktion kann selbst Dimensionswechsel auslösen, wenn gemischte Unterschiede
neue Achsen erzeugen, die den Kontinua isoliert nicht zur Verfügung stehen, und
wenn die Spannung die relevanten Schwellwerte überschreitet.

\subsection{Operatoralgebra und Nebenbedingungen}

Obwohl OC keine vollständige Operatoralgebra postuliert, sind mehrere
strukturelle Nebenbedingungen für innere Konsistenz erforderlich:

\paragraph{Kompatibilität mit Einbettungsräumen.}
Alle Operatoren müssen die Inklusion
\(A(K(t)) \subseteq A(M(t))\) respektieren und keine Entwicklung entlang von
Achsen erzeugen, die im \(M\) nicht vorhanden sind. Geburtsoperatoren sind nur
definiert, wenn \(M\) die relevanten Achsen bereits enthält.

\paragraph{Monotonie der Dimension.}
Der Strukturoperator \(S\) und die Geburtsoperatoren \(\Psi_{x\to x+1}\) müssen
erfüllen
\[
  \dim A(t+dt) \ge \dim A(t),
\]
mit strikter Ungleichheit nur bei Dimensionswechseln.

\paragraph{Schwellwertkonforme Dynamik.}
Operatoren \(F, G, H, Q, R\) müssen die durch Schwellwerte vorgegebene
Ungleichungsstruktur erhalten, außer bei autorisierten Überschreitungen von
kritischen oder Todesschwellwerten. Verletzungen von
\(\Theta_{\mathrm{exist}}\) werden als Todesereignisse interpretiert.

\paragraph{Continuumness als Viabilitätsindikator.}
Der Operator \(U\) koppelt alle anderen Komponenten; im Einzelnen:
\begin{itemize}
  \item wenn \(k(t) = 0\) und \(\Omega = \emptyset\), ist das Kontinuum tot und
        weitere Anwendungen von \(E\) haben auf dieser Ebene keine Wirkung;
  \item wenn \(k(t) > 0\), existiert mindestens ein stützender Zyklus und die
        Flüsse müssen die in \(F\) und \(Q\) kodierten Ausgleichsbedingungen
        erfüllen.
\end{itemize}

\paragraph{Kreuzebenen-Konsistenz.}
Für Übergänge \(K_x \to K_{x+1}\) muss die Operatorfamilie der höheren Ebene auf
die der niedrigeren Ebene zurückführen, wenn die neue Achse festgehalten und die
zusätzlichen Freiheitsgrade eingefroren werden. Damit sind die unteren Ebenen
als Spezialfälle höherer Ebenen wieder auffindbar, und die vertikale
Kohärenz der Hierarchie bleibt erhalten.

\subsection{Zusammenfassung}

Die Operatorfamilie \((F, G, H, Q, R, S, U, \Psi, E_{\mathrm{int}})\) stellt
das dynamische Rückgrat der Ontology of Continua dar. Statt einer einzelnen
hoche spezifischen Bewegungsgleichung verwendet OC ein modulares Satz von
Operatoren, die auf Flüsse, Potentiale, Schwellwerte, Zyklen, Grenzen, Achsen
und Continuumness wirken. Ihre kombinierte Wirkung regelt Geburt, Leben,
Interaktion und Tod von Kontinua über alle Ebenen \(K_0\)–\(K_{10}\).

Dieses Kapitel vervollständigt den formalen Kern von Core~1.3.3, indem es die
strukturelle Dynamik explizit macht, die in früheren Entwürfen nur implizit
vorkam, und den Boden für vollständig quantitative Modelle in nachfolgenden
Erweiterungspapieren bereitet.
